\section{Geometric and speed reductions}\label{sec:geometry}

\subsection{Strict records and regeneration words}

Put \(h(x)=(\ell/\|\ell\|_\infty)\cdot x\), so that a lattice step
has height increment of absolute value at most one, and some step has
increment exactly one. By hypothesis \(h(X_n)\to+\infty\) almost surely.
Let \(D_h\) be the event of never going below the starting height and put
\(p_h=P_0(D_h)\). An \emph{ordinary strict record} is a first strict
exceedance of the running height maximum, with time zero counted as the
initial record. Such a record is \emph{true} if the walk never subsequently
goes below its height.

We will repeatedly use a basic consequence of independence over sites.
The quenched weight of a finite path prefix uses only its departure rows.
A continuation event supported on paths avoiding those departure sites
has quenched probability measurable with respect to rows outside that set:
evaluate it using the walk killed on arrival at the forbidden set.
Its probability therefore factors from the prefix weight after annealing.
The same statement holds for infinite avoidance events by the cylinder
construction of path measures. For a finite continuation, it suffices
that its departure sites be outside the old departure set. None of these
facts requires independence of the entries within a row.

This construction belongs to the regeneration method of
Sznitman and Zerner~\cite{SznitmanZerner1999}; see also
\cite[Section~1]{Sznitman2002}. We give its needed
details, including arbitrary real height, directly.
The projected-width counting argument is related to the cone-regeneration
renewal argument in Simenhaus~\cite[Lemma~2]{Simenhaus2007}, who credits Zerner
for that renewal argument. Here we count ordinary record indices so that
the projected heights themselves need not form a lattice.

\begin{lemma}[True words and their widths]\label{geom-words}
We have \(p_h>0\). Under \(P_0(\,\cdot\mid D_h)\), include time zero as
a true boundary. The relative finite path words between consecutive true
boundaries are iid. The raw path has a finite first positive-time true
record, after which its translated suffix has this conditioned law.

For a conditioned word, let \(L>0\) be its height gain, let \(R\) be its
maximal Euclidean distance from its starting point, including its endpoint,
and let \(S_*\) count its ordinary strict records after its start and
through its end. Then
\[
0<\mathbb E L\le\mathbb E S_*\le p_h^{-1}.
\]
Here and until the change to coordinate height below, word expectations
refer to this conditioned word law.
\end{lemma}

\begin{proof}
The global height minimum is attained almost surely. Thus for some
deterministic \(n,x\) the event that \(X_n=x\) and that all later heights
are at least \(h(x)\) has positive probability. The quenched Markov
property at \(n\) gives
\[
E\big[P_{0,\omega}(X_n=x)P_{x,\omega}(D_h)\big]>0.
\]
Since the first factor is at most one, translation invariance gives
\(p_h>0\).

At a first strict record, every earlier departure row has smaller height.
The fresh-row rule therefore gives conditional raw probability \(p_h\)
that the record is true. Test the first positive-time strict record. If
an actual drop later detects its failure, test the next strict record
above the whole past observed by that failure time. All candidates and
failure detections are path stopping rules. Transience supplies another
candidate after every finite failure, and each candidate has success
probability \(p_h\). Thus a true record is eventually found almost surely.

For the iid assertion, enumerate the possible first words. A word from
\(0\) to \(z\) has positive length, every preterminal height in
\([0,h(z))\), and every intermediate strict record invalidated by a drop
before its end. Together with a suffix staying at or above \(h(z)\),
these conditions say exactly that this is the first word under \(D_h\).
Indeed a later suffix above \(h(z)\) cannot invalidate any earlier lower
record. The prefix weight and the suffix factor \(p_h\) use disjoint
departure rows. Dividing by the initial conditioning probability \(p_h\)
leaves the raw prefix weight as the word probability and an independent
translated suffix of the original conditioned law. Countable enumeration
and iteration give the iid words and infinitely many true boundaries.
For the first raw true record the same argument applies, without the
restriction that its prefix have nonnegative height. This argument does
not apply an annealed Markov property at a future-dependent cut.

Under \(D_h\), a fixed positive ordinary record index \(i\) is true with
probability equal to the raw probability of reaching that record before
dropping below zero: the fresh suffix contributes \(p_h\), which cancels
the conditioning. This probability is at least \(p_h\), since all record
indices are reached on \(D_h\). The true indices are a renewal process
with iid increments \(S_*\). Its positive-index count through \(m\),
divided by \(m\), converges almost surely and in expectation to
\(1/\mathbb E S_*\). This remains true with limit zero when the mean is
infinite, by the strong law on truncations; the ratios are bounded by one.
The preceding lower bound on each indicator therefore yields
\(\mathbb E S_*\le p_h^{-1}\). Each ordinary record increases height by
at most one, so \(L\le S_*\). This argument uses record indices, not a
lattice structure for the values of \(h\).
\end{proof}

\subsection{Finite spatial radius and the limiting ray}

The finite mean width controls progress only at true boundaries. To
control the path between them, we prove that the spatial radius of a
word also has finite mean. The argument compares two bounds for a
large displacement relative to the running height maximum. An infinite
mean radius would make such a displacement occur with probability of
order $1/B$. At each new ratio threshold, however, there is a fixed
positive chance of a fresh continuation that caps the ratio until the
walk leaves the height strip. Iterating these chances gives an
exponentially small upper bound in $B$.

\begin{proposition}[Integrable spatial radius]\label{geom-radius}
Under the conditioned word law,
\[
\mathbb E R<\infty. \tag{1}\label{eq:1}
\]
No moment of the word duration is assumed.
\end{proposition}

\begin{proof}
Suppose instead that $\mathbb E R=\infty$, and put
$I(s)=\mathbb E\min(R,s)$. This function is increasing, concave,
unbounded, and satisfies $I(s)=o(s)$, since $R$ is finite almost surely.
Write $H_k=\sum_{j=1}^kL_j$, with $H_0=0$.

\smallskip\noindent\emph{A family of likely finite templates.}
The width strong law gives constants $c>0$, $C$, and $C_0$ such that
\[
ck-C_0\le H_k\le Ck+C_0\qquad(k\ge0)
\]
holds with probability greater than $.95$. Fix an integer $C_1$ with
$C_1c>C+4$, then $b>1/c$, and then $\eta>0$ so small that
$4C_1\eta<.1$. For large $a$, choose the least positive integer $N$
with $I(aN)\ge\eta a$. Then $N\to\infty$, and minimality and
concavity give
\[
\eta a\le I(aN)<2\eta a
\]
for all sufficiently large $a$.

With probability greater than $.8$, the first $C_1N$ words satisfy the
height bounds above and
\[
\sum_{j=1}^kR_j\le ak\qquad(1\le k\le C_1N).
\tag{2}\label{eq:2}
\]
Here is the maximal estimate behind this assertion. Truncate each
radius at $aC_1N$. Its expectation is at most
$C_1I(aN)<2C_1\eta a$. The union bound for a truncation among the
first $C_1N$ radii is therefore at most $2C_1\eta$. For a stationary
nonnegative sequence $(Z_j)$,
\[
\Pr\!\left(\max_{1\le k\le n}\frac1k\sum_{j=1}^kZ_j>a\right)
\le \frac{\mathbb E Z_1}{a}.
\]
To see this, mark the starts among $1,\ldots,M$ admitting a witness
interval of length at most $n$ with average greater than $a$. Select
such intervals greedily from left to right. They are disjoint and cover
all marked starts, so their total sum is greater than $a$ times the
number of marked starts. The selected intervals lie in $1,\ldots,M+n$.
Take expectations, divide by $M$, and let $M\to\infty$ to remove
the boundary contribution.
Applied to the truncated radii, this costs at most another
$2C_1\eta$. Together with the height event this proves the asserted
$.8$ bound. This is the nonnegative-sequence form of Hopf's maximal
ergodic inequality; see Garsia~\cite{Garsia1965}.

\smallskip\noindent\emph{The polynomial lower bound.}
Fix a large $B>1$. Choose $i_0$ so that
\[
\sum_{i\ge i_0}\Pr(L>i)<\frac{\eta}{4B},\qquad
C(i-1)+C_0+i\le(C+2)i\quad(i\ge i_0).
\]
Choose $a$ sufficiently large that $N\ge i_0$ and
$I(Bai_0)<\eta a/4$. Tail integration, followed by removal of the
width exceptions, gives
\[
\frac{\eta}{2B}\le
\lambda:=\sum_{i=i_0}^N\Pr(R>Bai,\ L\le i)
\le2\eta.
\]
For the lower bound, the unrestricted sum is at least
$[I(Ba(N+1))-I(Bai_0)]/(Ba)\ge3\eta/(4B)$; for the upper bound it
is at most $I(BaN)/(Ba)\le I(aN)/a$.

Count the indices $i$ at which this large-word event occurs and the
preceding words satisfy the height and radius bounds through $i-1$.
Independence gives mean at least $.8\lambda$. Dropping the
preceding-word restrictions bounds the second moment by
$\lambda+\lambda^2$. The count is therefore positive with probability
at least $c_2/B$, where $c_2>0$ does not depend on $B$ or $a$.
At some point of that word,
\[
\|X_n\|>(B-1)ai,
\qquad 0\le h(X_j),\quad
M_n:=\max_{j\le n}h(X_j)\le(C+2)i.
\]
The height bound holds throughout this prefix because every word ends
at a strict record and has no earlier height above its endpoint.
One of the $2d$ signed coordinate vectors $u$ thus satisfies
\[
\frac{Y(X_n)}{a(1+M_n)}>c_3B,
\qquad Y(x)=u\cdot x+ba h(x),
\tag{3}\label{eq:3}
\]
for a fixed $c_3>0$ and sufficiently large fixed $B$.
Consequently, under the raw law the union of these events over $u$
has probability at least $p_hc_2/B$, before the path leaves
$\{h\ge0,\ M\le(C+2)N\}$.

\smallskip\noindent\emph{A fresh continuation caps the ratio.}
Fix $u$ and put $Q_n=Y(X_n)/[a(1+M_n)]$. Let $\zeta$ be the exit
time just described. We claim that constants $\epsilon_0>0$ and
$K<\infty$, independent of $B,a$ and the threshold $t>0$, have the
following property: at every finite prefix ending at the first crossing
$\sigma<\zeta$ of $Q>t$, a continuation of conditional raw probability
at least $\epsilon_0$ exits the strip before $Q$ exceeds $t+K$.

At the crossing,
\[
Y(X_j)\le ta(1+M_j)\le ta(1+M_\sigma)<Y(X_\sigma)
\quad(j<\sigma).
\]
Thus the arrival is a strict $Y$-record. A step changes $Y$ by at most
$ba+1$, so $Q_\sigma\le t+b+1$ when $a\ge1$.
Choose an integer $M_0'>C_0$ with
$M_0'b>1+bC_0$. First prescribe $M_0'$ steps of height $+1$.
Each increases $Y$ by at least $ba-1>0$. At their endpoint impose a
translated $D_h$-suffix whose first $C_1N$ words satisfy the template
above. In its word $k+1$, every site has $Y$-increment from the launch
at least
\[
M_0'(ba-1)+ba(ck-C_0)-a(k+1)>0.
\]
This holds uniformly in $0\le k<C_1N$ for large $a$: the coefficient
of $k$ is $a(bc-1)>0$, and the constant term is positive by the choice
of $M_0'$. These words therefore avoid every old departure row.
Their terminal true cut has height at least
\[
h(X_\sigma)+M_0'+cC_1N-C_0>(C+2)N.
\]
The entire remaining suffix stays above this cut, and so also avoids
all old departures. The template stays above its own starting height,
which is strictly above every script departure. Thus the whole infinite
continuation uses no old departure row: first the $Y$ barrier provides
separation, and then the terminal height does. The fresh-row
factorization gives probability at least
\[
\epsilon_0=.8\,\kappa^{M_0'}p_h>0.
\]
Figure~\ref{fig:forward-template} records these two separation mechanisms.

\begin{figure}[tbp]
\centering
\begin{tikzpicture}[x=1cm,y=.83cm,font=\small,>=Stealth,
  old/.style={draw=black!55,line width=1pt},
  script/.style={draw=orange!80!black,line width=1.6pt},
  controlled/.style={draw=blue!65!black,line width=1.3pt},
  tail/.style={draw=green!45!black,line width=1.3pt}]
  \fill[black!5] (-4.3,.2) rectangle (0,4);
  \fill[green!5] (-4.3,5.15) rectangle (7,6.5);
  \draw[->,black!65] (-4.45,.05) -- (7.25,.05)
    node[below left] {$Y-Y(X_{\rm launch})$};
  \draw[->,black!65] (-4.45,.05) -- (-4.45,6.7) node[above] {$h$};
  \draw[densely dashed,black!55] (0,.15) -- (0,4.85);
  \draw[densely dashed,black!45] (-4.3,4) -- (6.9,4);
  \node[anchor=south west,fill=white,inner sep=2pt] at (-4.1,4)
    {old height maximum};
  \draw[old] (-3.9,1.15) -- (-3.5,1.45) -- (-3.8,1.6)
    -- (-2.9,1.85) -- (-3.15,1.95) -- (-2.1,4)
    -- (-1.5,2.7) -- (-.9,2.1) -- (-1.25,2.4) -- (0,2.2);
  \node[align=center,anchor=south] at (-2.15,.2)
    {old departure\\rows have $Y<Y(X_{\rm launch})$};
  \fill (0,2.2) circle (2pt);
  \draw[black!55] (.15,2.1) -- (1.2,1.4);
  \node[anchor=west,align=left] at (1.3,1.2)
    {launching\\strict $Y$-record};
  \draw[script,->] (0,2.2) -- (.85,3.45);
  \node[anchor=west,text=orange!80!black,align=left] at (1,2.2)
    {$M_0'$ upward steps};
  \draw[densely dotted,orange!75!black] (.85,3.45) -- (6.8,3.45);
  \node[anchor=north east,text=orange!75!black] at (6.8,3.45)
    {template base};
  \draw[controlled,->] (.85,3.45) -- (1.35,3.85) -- (1.05,3.65)
    -- (2.15,4.05) -- (1.8,3.85) -- (2.9,4.4)
    -- (2.55,4.15) -- (3.65,4.65) -- (3.2,4.4) -- (4.15,5.15);
  \node[anchor=west,align=left,text=blue!65!black] at (4.3,4.52)
    {first $C_1N$\\controlled words};
  \fill[blue!65!black] (4.15,5.15) circle (2pt);
  \draw[densely dashed,green!45!black] (-4.3,5.15) -- (7,5.15);
  \node[anchor=south west,fill=white,inner sep=2pt] at (-4.1,5.15)
    {true cut above all old heights};
  \draw[tail,->] (4.15,5.15) -- (4.65,5.6) -- (3.5,5.8)
    -- (2.25,5.4) -- (1.15,6.05) -- (-.65,5.8) -- (-1.65,6.35);
  \node[anchor=west,text=green!45!black,align=left] at (4.8,5.85)
    {infinite $D_h$\\continuation};
\end{tikzpicture}

\caption{The forward continuation in the $(Y,h)$ projection
(schematic, not to scale). Before launch, all departure rows have smaller
$Y$. The upward script raises the template base. The controlled words
stay beyond the old $Y$ maximum and above the script departures. Their
terminal true cut lies above all old heights; after that cut the suffix
may cross the old $Y$ barrier while still avoiding the old rows. The
projection does not restrict the walk to two dimensions.}
\label{fig:forward-template}
\end{figure}

It remains to bound the ratio on this continuation. The positive part
of the launch coordinate $u\cdot X_\sigma$, divided by any later
denominator, is at most $t+b+1$, since $h(X_\sigma)\ge0$ and the
denominator is nondecreasing. The later height term contributes at most
$b$. The script adds at most $M_0'$ to the ratio. In word $k+1$, the
additional coordinate contribution is at most
\[
\frac{M_0'/a+k+1}{A_0+ck}
\le\max\left\{\frac{M_0'+1}{A_0},\frac1c\right\},
\qquad A_0=1+M_0'-C_0>1.
\]
The numerator uses the bound on the sum of radii through word $k+1$;
the denominator uses the height of that word's starting boundary.
Thus any fixed
\[
K>2b+1+\max\left\{M_0',\frac{M_0'+1}{A_0},\frac1c\right\}
\]
proves the claim. Exit occurs by the last prescribed template endpoint.

\smallskip\noindent\emph{The exponential upper bound.}
Let $\sigma_j$ be the first crossing before $\zeta$ of
$t_j=1+jK$. The continuation just constructed prevents
$\sigma_{j+1}<\zeta$. Summing its bound over the finite prefixes at
$\sigma_j$ yields
\[
P_0(\sigma_{j+1}<\zeta)
\le(1-\epsilon_0)P_0(\sigma_j<\zeta).
\]
This conditions only on the observed crossing prefix; it does not reveal
whether an infinite template has failed. Hence the event in
\eqref{eq:3} has probability at most $C'e^{-c'B}$ for this $u$.
Taking the union over signed coordinates contradicts the lower bound
$p_hc_2/B$ when $B$ is large. All estimates use the permitted order:
fix the geometry constants, then $B$, then $i_0$, then take $a$ large.
This proves \eqref{eq:1}.
\end{proof}

\begin{corollary}[Deterministic limiting ray]\label{geom-ray}
There is a deterministic unit vector \(r\) with \(h(r)>0\) such that
\(X_n/\|X_n\|\to r\) almost surely. After a deterministic signed
permutation of the coordinate axes, \(r_1>0\), and all the preceding
word conclusions hold with height \(x_1\).
\end{corollary}

\begin{proof}
Word endpoints obey the vector strong law, with mean displacement
\(w\) satisfying \(h(w)=\mathbb E L>0\). Proposition~\ref{geom-radius}
also gives \(R_k/k\to0\) almost surely: for each \(\varepsilon>0\),
the probabilities \(\Pr(R_k>\varepsilon k)\) are summable. Thus
interpolation inside each word gives the same limiting direction
\(r=w/\|w\|\), without any duration moment. The completed-word count
tends to infinity. Lemma~\ref{geom-words} transfers the conclusion
through the finite initial raw prefix.

Some signed coordinate of \(r\) is positive. Relabel it as the first
coordinate, preserving the lattice, uniform ellipticity, and the iid
row law. Then \((X_n)_1\to+\infty\) almost surely. The proofs of
Lemma~\ref{geom-words} and Proposition~\ref{geom-radius} apply again
with this coordinate height, using only its transience. We retain the
same notation for the transformed walks and laws.
\end{proof}

\subsection{Layer notation; a small inverse moment}

Henceforth height means \(x_1\). Write \(x=(x_1,\underline x)\), \(\underline x\in\mathbb Z^{d-1}\); for fluctuations we use the single scalar coordinate \(x_2\). For a path from \(x\) and \(j\in\mathbb Z\), \(T_j=\inf\{n\ge0:(X_n)_1=x_1+j\}\) denotes first hit of the layer \(x_1+j\), relative to the indicated start (\(\inf\varnothing=\infty\); \(T_0=0\)). Put \(D=\{T_{-1}=\infty\}\), \(p=P_x(D)>0\),
\[
P_x^+=P_x(\,\cdot\,\mid D),\qquad q(x)=P_{x,\omega}(D),\qquad
a_H(x)=P_{x,\omega}(T_H<T_{-1})\quad(H\ge0\ \text{integer}),
\]
where the environment in notation such as \(q,a\) is implicit. Let \(K_H(x,\cdot)\) be the unnormalized quenched kernel of reaching that upper layer before the drop, of mass \(a_H(x)\); we also use it for the whole path prefix through this first hit, testing events on that prefix. It only uses rows in the \(H\) layers from the starting layer inclusive to the upper one exclusive. From 0 abbreviate \(a_H=a_H(0)\), \(u_H=E a_H\). The filtration \(\mathcal F_H\) on environments reveals rows at \(0\le x_1<H\); under \(P\) the upper layers from \(H\) inclusive are thus fresh iid layers, also starting at any finite \(\mathcal F_H\)-stopping layer (by partitioning according to that layer). Corresponding shifted filtrations at other bases will be used similarly. Write \(L_{\rm s},R_{\rm s},\tau_{\rm s}\) for the height width, radius and number of path steps in a single-walk true word in direction \(e_1\) under \(P_0^+\).

\begin{lemma}[A small inverse moment]\label{geom-inverse}
For every site \(x\), \(a_H(x)\downarrow q(x)>0\) for almost every
environment. Moreover, for some \(\lambda>0\),
\[
E\,q(0)^{-\lambda}<\infty. \tag{4}\label{eq:4}
\]
\end{lemma}

\begin{proof}
Uniform ellipticity excludes confinement forever in a fixed finite-height
strip, even quenched: blocks of sufficiently many consecutive upward
steps have a fixed positive chance. Hence \(a_H(x)\downarrow q(x)\).
Also \(q(je_1)\) uses only layers at or above \(j\), so the event that
\(q(je_1)>0\) infinitely often is a tail event of the independent layers.
For every \(m\), the union of these events over \(j\ge m\) has
probability at least \(P(q(0)>0)>0\), by stationarity and \(Eq(0)=p>0\).
The tail zero--one law makes their limsup almost sure. A finite upward
script connects zero to one such site and gives \(q(0)>0\). Translation
and countability give positivity at every site simultaneously.

At layer \(H\), let \(\mu_H=K_H(0,\cdot)/a_H\) be the endpoint
probability distribution. It is \(\mathcal F_H\)-measurable, and the
quenched Markov property, restricted to a suffix never below layer \(H\),
gives
\[
q(0)\ \ge\ a_H\int q(x)\,\mu_H(dx).
\]
Each \(q(x)\) on that layer has the original marginal law independently
of the lower rows. If \(q(0)/a_H<s\), then more than half the
\(\mu_H\)-mass has \(q(x)\le2s\). Conditional expectation followed by
Markov's inequality therefore gives
\[
P\big(q(0)/a_H<s\mid\mathcal F_H\big)
\le2P\big(q(0)\le2s\big).
\]
The endpoint values \(q(x)\) need not be independent of one another.
Partitioning by a finite stopping layer gives the same conditional bound
there.

Choose \(0<\delta<\kappa\) sufficiently small that
\(\rho:=2P(q(0)\le2\delta/\kappa)<1\), and set
\[
H_j=\inf\{H\ge0:a_H<\delta^j\},\qquad j\ge1.
\]
These are stopping layers. One upward step after a successful crossing
gives \(a_{H+1}\ge\kappa a_H\), so on \(H_j<\infty\),
\(a_{H_j}\ge\kappa\delta^j\). The event \(H_{j+1}<\infty\) implies
\(q(0)<\delta^{j+1}\), and hence
\(q(0)/a_{H_j}<\delta/\kappa\). The preceding conditional bound yields
\[
P(H_{j+1}<\infty\mid\mathcal F_{H_j})\le\rho
\quad\text{on }\{H_j<\infty\}.
\]
Since \(a_H\downarrow q(0)\), the event \(H_j<\infty\) is exactly
\(\{q(0)<\delta^j\}\). Iteration gives its probability at most
\(\rho^{j-1}\). Summing over
\(\delta^{j+1}\le q(0)<\delta^j\), and choosing \(\lambda>0\) with
\(\rho\delta^{-\lambda}<1\), proves \eqref{eq:4}.
\end{proof}

\subsection{Probability estimate sufficient for speed}

The rest of the paper will establish the following estimate for one
fixed \(0<\beta<1\):
\[
P(a_N<N^{-\beta})<N^{-D_0}
\quad\text{eventually for each fixed finite }D_0>0.
\tag{5}\label{eq:5}
\]

Small quenched exit probabilities also connect spatial estimates to
regeneration-time tails in Sznitman's
work~\cite[Lemma~3.2 and Theorem~3.3]{Sznitman2002}.
The next proposition gives the reduction for the crossing probabilities
defined here, with \eqref{eq:5} as its explicit input.

\begin{proposition}[From crossing probabilities to speed]\label{geom-speed}
If \eqref{eq:5} holds, then \(\mathbb E\tau_{\rm s}<\infty\), and
the walk has an almost-sure deterministic velocity with strictly positive
projection on the original direction \(\ell\).
\end{proposition}

\begin{proof}
All conclusions below using \eqref{eq:5} are conditional on that estimate.

\smallskip\noindent\emph{Spatial confinement before exit.}
The raw path before leaving the height interval \([0,N]\) stays in
\(\|\underline X\|_\infty\le CN^2\), except with probability at most
\(Ce^{-cN}\). To prove this, fix a signed lateral coordinate \(u\).
By Corollary~\ref{geom-ray}, a sufficiently large fixed \(b>0\) makes
\(Y(x)=u\cdot x+bx_1\) transient in its positive direction. Thus the
event of staying forever at \(Y\ge0\) has positive raw probability,
by the minimum argument of Lemma~\ref{geom-words}.

For a sufficiently large fixed \(C_2\), the path through its first
relative height hit \(N+1\) has norm at most \(C_2N\) with probability
tending to one. Indeed the ray gives
\(\|X_n\|\le(C_2/2)(X_n)_1\) eventually, and the earlier path has a
finite maximum norm. The intersection of this event with the preceding
no-backtracking event therefore has probability bounded below uniformly
for large \(N\).

At a first crossing of \(Y>jC_3N\), before strip exit, the walk is at
a strict \(Y\)-record. Impose that favorable translated suffix using
the fresh rows at or above the record. If \(C_3\) is sufficiently large,
the suffix precludes crossing \((j+1)C_3N\) before strip exit: its
\(Y\)-displacement through relative height \(N+1\) is at most
\((1+b)C_2N\), the launch overshoot is bounded, and absolute strip
exit occurs no later than this relative height hit. The finite-prefix
threshold recursion used in Proposition~\ref{geom-radius}, now with
spacing \(C_3N\), gives exponential decay through \(N\) thresholds.
Since \(u\cdot X\le Y\) before exit, the finite union over signed
lateral coordinates proves the confinement assertion.

\smallskip\noindent\emph{Super-polynomial width tails.}

Now \(\Delta u_N:=u_N-u_{N+1}\ge0\) is the annealed probability of reaching \(N\) before the drop but dropping before \(N+1\). In environments good at all sites of the strip and box (where \(a_{N+1}(x)\ge(N+1)^{-\beta}\)), the portion of this event inside the box has quenched probability at most
\[
(1-(N+1)^{-\beta})^{N+1}.
\]
Indeed, after reaching \(N\), at first descending-level visits to each of \(N,N-1,\ldots,0\) still in the box, there is at least the indicated chance to hit top \(N+1\) before the next descent by one. Condition \eqref{eq:5}, translation, a polynomial union bound, and the box claim thus give super-polynomial decay of \(\Delta u_N\).

The true levels under \(P_0^+\) form the renewal process of iid widths \(L_{\rm s}\), with renewal masses exactly \(u_n\): impose the fresh \(D\)-suffix at the first reach of \(n\) before drop, and cancel \(p\). If \(G_j=P_0^+(L_{\rm s}>j)\), then \(\sum_{j=0}^n G_j u_{n-j}=1\) by the last renewal up to \(n\), so
\[
G_{n+1}=\sum_{j=0}^n G_j\Delta u_{n-j}.
\]
Put \(f_k=\Delta u_{k-1}\) for \(k\ge1\). This nonnegative kernel has
total mass \(\theta=1-p<1\) and all polynomial moments. The recurrence,
with \(G_0=1\), gives \(G=\sum_{r\ge0}f^{*r}\), including unit mass
at zero for \(r=0\). If \(\theta>0\), normalize \(f\) by \(\theta\).
For every positive integer \(m\), the \(m\)-th moment of its
\(r\)-fold convolution is at most \(r^m\) times its \(m\)-th moment,
by
\((z_1+\cdots+z_r)^m\le r^{m-1}\sum_i z_i^m\).
Summing with the geometric weights \(\theta^r\) proves
\(\sum_n n^mG_n<\infty\). The case \(\theta=0\) is immediate.
In particular the width tail \(G_n\) decays super-polynomially.

\smallskip\noindent\emph{Time integrability and velocity.}
The raw exit time \(T_N\wedge T_{-1}\) exceeds \(N^A\) with
super-polynomially small probability, for a sufficiently large fixed
\(A\). To see this, use confinement and \eqref{eq:5} on a slightly
enlarged box and strip, requiring
\(a_N(x+e_1)\ge N^{-\beta}\) for every relevant site with
\(0\le x_1<N\). A polynomial union bound controls the exceptional
environments.

From any visit to such a site \(x\), the quenched probability of no
further visit before exit is at least \(\kappa a_N(x+e_1)\): take one
upward step and then advance \(N\) before dropping below the new layer.
This event forces original strip exit before revisiting \(x\), even
though the auxiliary target may be higher than \(N\). The quenched
Markov property consequently bounds each site's visit count by a
geometric tail with success probability at least \(\kappa N^{-\beta}\).
There are at most \(CN^{2d-1}\) sites in the confined box. If the exit
time exceeds \(N^A\), one site has at least
\(N^A/(CN^{2d-1})\) visits. For \(A>2d-1+\beta\), the union of these
geometric-tail events is super-polynomially small, as are the bad-box
and bad-environment events.

On \(D\), if \(L_{\rm s}\le N\), the first word ends at its first
arrival to height \(L_{\rm s}\), and hence
\(\tau_{\rm s}\le T_N\). Thus
\[
P_0^+(\tau_{\rm s}>N^A)
\le P_0^+(L_{\rm s}>N)
+p^{-1}P_0(T_N\wedge T_{-1}>N^A).
\]
Both terms have super-polynomial decay, proving
\(\mathbb E\tau_{\rm s}<\infty\). The iid-word strong laws now give
the velocity as mean displacement divided by mean duration; its first
coordinate is \(\mathbb E L_{\rm s}/\mathbb E\tau_{\rm s}>0\).
The integrable spatial radius controls the inter-endpoint deviations,
so this limit holds at every path time. The finite raw initial prefix
does not change it. Finally the velocity is a positive multiple of the
ray in Corollary~\ref{geom-ray}, which has positive projection on the
original \(\ell\). Undoing the signed coordinate permutation proves
the stated conclusion. As an almost-sure deterministic limit, this
velocity is unique for the law \(\nu\).
\end{proof}

It remains to prove \eqref{eq:5}. Sections~\ref{sec:fluctuations}--\ref{sec:stationary}
are devoted to that estimate; no finite duration moment is used there.
